\documentclass[../main.tex]{subfiles}
\begin{document}
%==========================================TIÊU ĐỀ TẬP========================================
\begin{table}[h]
  \begin{adjustwidth}{-1cm}{}
    \begin{tabular}{>{\centering\arraybackslash}p{6cm}|>{\raggedright\arraybackslash}p{12.5cm}}
      \multirow{5}{6cm}
      % Cột bên trái: ÉPISODE và số tập
      {\\[-40pt]
        \flaregothic\fontsize{20pt}{20pt}\selectfont \centering
        \color{eptype}ÉPISODE\color{black}\\
        \burbank\fontsize{72pt}{72pt}\selectfont
        \color{epnum}106
        \\[18pt]
        \flaregothic\fontsize{14pt}{14pt}\selectfont
        \color{epattr}BIENVENUE, \uline{\color{Neneseal} Nene} !\\
        \color{black}
      }
       &
      % Dòng thứ nhất: tên tập tiếng Nhật
      {\fotskip\fontsize{18pt}{18pt}\selectfont \ruby{春}{はる}はあげぽよ〜＾＾v
      }                                                                   \\[6pt]
      % Dòng thứ hai: tên tập tiếng Anh
       & {\cafeteria\fontsize{18pt}{18pt}\selectfont In Spring, w00t!}                                                                 \\[4pt]
      % Dòng thứ ba: tên tập tiếng Việt
       & {\vnmsans\fontsize{18pt}{18pt}\selectfont Mùa xuân, mùa Nene}                                                                 \\[8pt]
      % Dòng thứ tư: Nhân vật xuất hiện
       & {\sen{\fontsize{14pt}{14pt}\selectfont Track used: Gemini (Luigi's Cool Pants - F-777) (\ruby{二}{に}\ruby{重}{じゅう}\ruby{星}{せい})}
         }
      \\[8pt]
      % Dòng cuối cùng: Thời gian diễn ra của tập (thường sẽ là ngày phát hành)
       & {\continuum\fontsize{16pt}{16pt}\selectfont Ngày phát hành: 16/05/2021}                                                       \\[18pt]
    \end{tabular}
  \end{adjustwidth}
\end{table}
%=========================================TRUYỆN CHÍNH========================================
\color{black}

\pov{Hikawa Kuon}

Đã ba tháng rồi kể từ thành viên thứ ba của \textit{NePoLaBo} lên văn phòng \hll\ lần đầu (và thứ tư, nếu tính Aloe-chan, nhưng ở nhà tôi.)

Tôi đã được Botan-chan bảo rằng là thành viên cuối cùng của Thế hệ thứ Năm sẽ lên đây.

Một idol đến từ vũ trụ khác (tôi cho là vậy), tràn đầy năng lượng và luôn mang theo mình những ước mơ rực rỡ. Nếu Botan là một sư tử trầm ổn, thì người này lại là một mặt trời nhỏ, lúc nào cũng bùng nổ sự vui tươi và chẳng bao giờ để bầu không khí trở nên u ám.

Em ấy sở hữu một sức hút đầy tự nhiên: dù có lúc hậu đậu, có lúc ngốc nghếch, nhưng sự hồn nhiên và chân thành của em ấy lại khiến ai ở bên cạnh cũng cảm thấy vui vẻ. Với giọng hát trong trẻo và tinh thần làm việc không ngừng nghỉ, em ấy luôn khao khát vươn lên để chạm tới giấc mơ trở thành idol số một.

Bên cạnh đó, tính cách của em cũng rất đáng yêu: ngây thơ, nhí nhố khi ở bên những người thân thiết như Polka hay Lamy, nhưng cũng có những khoảnh khắc trưởng thành và nghiêm túc đến bất ngờ. Dù bị trêu chọc là ``một con hải cẩu ngốc nghếch'', nhưng thực chất em ấy rất nỗ lực để hoàn thiện bản thân mỗi ngày. Một cô gái không ngại sai sót, chỉ cần có thể khiến mọi người mỉm cười.

Ngoài ra, em ấy còn là một người rất hay\dots\ lập một bộ sưu tập chồng, bất kể người đó tới từ nơi nào hay bao nhiêu tuổi.

Và dù là một thần tựng, mục tiêu của em không chỉ đơn giản là hát và nhảy. Em ấy muốn một ngày nào đó có thể đứng trên sân khấu lớn nhất, để tất cả mọi người đều biết đến mình, không chỉ vì sự vui nhộn mà còn vì những nỗ lực không ngừng nghỉ của mình.

Một idol với giấc mơ vươn xa.

\dots Hoặc có lẽ, một ``vị thần'' ngố tự phong đang trên đường chinh phục thế giới.

\chrint

Người mà tôi đang nói đến là Momosuzu Nene (\ruby{桃}{もも}\ruby{鈴}{すず}ねね, \ruby{Mô-mô}{Đào}-\ruby{xư-dư}{Linh}\ Nê-nê) - thề luôn, cái tên này tôi đã đọc sai không biết bao nhiêu lần - là thành viên thứ tư (hoặc thứ năm) của thế hệ thứ Năm. Một cô gái có ba tài năng: họa-ca-múa, và cô nàng chau chuốt sở trường của mình một cách chăm chỉ mỗi ngày.

\begin{wrapfigure}[25]{l}{8.5cm}
  \includegraphics[width=8.5cm]{../../../../Image/Book?/nene_model.png}
\end{wrapfigure}

Nene-chan là một cô nàng vui vẻ, hoạt bát và thân thiện, thích nói chuyện và ca hát, có một mong muốn lớn lao là trở thành một thần tượng và có bài hát mở đầu anime của riêng mình.

Tính cách trong trẻo và dễ mến ấy còn được thể hiện ở việc thường tỏ ra tự mãn dễ thương khi thắng trò chơi và bĩu môi mỗi khi thua. Em ấy cũng là một ``kẻ đa tình'' khi có rất nhiều ``chồng'' - thực chất là những người xem tỏ vẻ trung thành với em ấy, dù là nam hay nữ - nhưng đồng thời còn tán tỉnh các thành viên \hll\ khác một cách vui vẻ mà em ấy gọi là ``vợ'', với một số ví dụ bao gồm Lamy-chan và Ina-ina\footnote{Sẽ được giới thiệu ở quyển Mười Ba, chính thức xuất hiện ở quyển Mười Sáu.}.

Cô nàng ``Hải cẩu'' này được Pol-pol và Lamy-chan nhận xét là thành viên có tính cách mạnh mẽ nhất trong nhóm, vì em ấy có thể dễ dàng chuyển về tâm trạng vui tươi và lạc quan. Tôi nghe nói em ấy đã tự nộp đơn gia nhập ba lần trước đó, trước khi thật sự thành công vào lần thứ tư.

Ngoài ra, dường như em ấy cũng là người có cảm xúc thuộc vào dạng mãnh liệt nhất trong Thế hệ, thể hiện ở việc nhớ nhung thành viên thất lạc của mình, tới mức bạo bệnh tận một tháng trời - theo những gì mà ba thành viên còn lại kể cho tôi.

Một điểm nữa mà chỉ có Nene-chi có, là em ấy hay phát trực tiếp vào giờ ăn trưa để trò chuyện với người hâm mộ và hay trả lời những câu hỏi dí dỏm và hài hước từ những người xem của mình.

Thỉnh thoảng, em ấy sẽ nói một vài câu cơ bản bằng tiếng Anh và tiếng Tây Ban Nha khi tương tác với khán giả, dẫn đến số lượng người hâm mộ nước ngoài tăng lên trên các buổi phát trực tiếp của mình (dù bản thân không thạo ngoại ngữ lắm). Vì lẽ đó nên người hâm mộ đã đặt cho em biệt danh ``Worldwide Nenechi (Nene-chi toàn cầu)'' bất cứ khi nào điều này xảy ra.

Nene-chan cũng có vẻ như có biệt danh mới mỗi lần em ấy phát trực tiếp, với một cái tên cực dài mà khi tôi đọc xong cũng méo cả mồm: ``Super Hyper Ultra Ultimate Deluxe Perfect Amazing Shining God 東方不敗 (Đông Phương Bất Bại) Master Ginga Victory Strong Cute Beautiful Galaxy Baby 無限 無敵 無双 (Vô Hạn, Vô Địch, Vô Song) NENECHI!''

Nene-chan có một đôi mắt xanh lá cây, một mái tóc dài màu vàng được chuyển dần sang hồng khi xuống dưới và được búi hai bên, được trang trí bằng dây buộc hoa, phần mái cũng được rẽ ngôi đôi.

Bộ trang phục của cô nàng thường ngày bao gồm: một chiếc áo sơ mi không tay màu cam có diềm xếp nếp với những đường cắt ở khe ngực và phía trên rốn, một chiếc áo nịt dưới ngực màu nâu có đan tréo, hai tay áo ngắn rời, một chiếc nơ mang ba màu vàng-cam-nâu và đeo chiếc chuông nhỏ hình quả đào ở cổ; một bộ đầm ngắn xếp nếp màu cam và vàng được chia thành nhiều lớp, một chiếc váy lót màu trắng áng hồng có hoa trên đó, một chiếc nơ lớn buộc sau lưng, một chiếc túi xách hình gấu (thường cô nàng sẽ không đeo), bên ngoài có mặc một chiếc áo khoác thể thao (dùng cho sinh viên/học sinh) màu trắng-cam. Nene-chi đi một chiếc tất có diềm xếp nếp dài đến đùi và một chiếc tất ngắn hơn buộc nơ vàng, đi đôi bốt màu nâu.

Hình như tôi chưa nói là búi tóc của em ấy có thể biến thành tai gấu (giả) và có cả bộ găng tay gấu không nhỉ?

\chrint

Không biết là cô nàng Gấu-Hải cẩu này sẽ có trò gì cho chúng tôi đây\dots

\il

Một ngày trung tuần tháng Năm. Mùa xuân đã trôi qua một nửa, nhưng mùa hoa đang rực rỡ nhất. Những hàng cây anh đào dọc hai bên đường phủ đầy sắc hồng phấn, thi thoảng có cơn gió nhẹ lướt qua cuốn theo những cánh hoa bay lả tả trong không trung. Ánh nắng xuân dịu dàng chiếu xuống, khiến khung cảnh càng thêm phần thơ mộng.

Tôi thì không vấn đề gì, vì tôi không bị dị ứng với phấn hoa. Nhưng các cô gái đi cùng tôi, đặc biệt là thành viên mới đây, thì lại không may mắn như vậy.

\begin{dialogue}
  \speak{Kanata} Hắt xì!
  \speak{Flare} Xì!
\end{dialogue}

\begin{figure}[H]
  \centering
  \begin{subfigure}[t]{0.45\textwidth}
    \centering
    \includegraphics[width=8cm]{../../../../Image/Book?/8-10a1.png}
  \end{subfigure}
  \quad
  \begin{subfigure}[t]{0.45\textwidth}
    \centering
    \includegraphics[width=8cm]{../../../../Image/Book?/8-10a2.png}
  \end{subfigure}
  \caption{Cách mà Kanata và Flare hắt hơi.}
\end{figure}

Kanata vừa hắt hơi một cái rõ to, Flare cũng khẽ xoa mũi, còn Momosuzu Nene - người vừa đa tài lại vừa\dots\ có phần quái đản - lại hắt hơi theo một phong cách không ai ngờ tới: ``\textbf{Momosu-dzu Nene!}''

\img{8-10a3}

Đúng. Là hắt xì một cách RẤT cụ thể. Thay vì là một tiếng ``ắt xì!'' như bình thường.

Tôi chớp mắt, Kanata cũng giật mình nhìn cô ấy với vẻ khó hiểu. Tôi nghiêng đầu nhìn chủ nhân vừa có cú hắt hơi kỳ quặc đó: ``Anh chưa thấy ai hắt xì dị như em đấy.''

``Cái đó mà là hắt hơi à?'' Kana-chan cũng tiếp lời, thắc mắc với kiểu hắt hơi dị ứng của Nene-chi.

Thế nhưng cô nàng ấy lại không hề cảm thấy có gì kỳ lạ. Trái lại, em ấy còn giơ hai ngón tay cái lên như thể rất tự hào về điều đó. Tôi chỉ có thể thở dài: ``Đó không phải là lời khen đâu\dots''

Khi chúng tôi còn đang bối rối vì kiểu hắt hơi kỳ quặc của em ấy, Flare-chan đột nhiên lên tiếng, chuyển chủ đề: ``Mọi người có biết là mùa xuân bắt đầu lúc tiết lập xuân (立春, bắt đầu mùa xuân) không?''

Hai người con gái kia có vẻ chưa từng nghe đến khái niệm này bao giờ, còn tôi thì thẳng thừng đáp: ``Tầm khoảng ngày 20 hay 21 tháng Ba là vào xuân chứ? Có, anh từng đọc qua rồi.''

Em ấy khẽ gật đầu như hài lòng với câu trả lời của tôi, nhưng chưa kịp nói thêm gì thì Nene đột nhiên ngơ ngác hỏi lại: ``Điểm\dots\ thân phận\footnote{Nguyên văn: 立直, nghĩa của nó là tên một tay bài trong mạt chược kiểu Nhật.}?''

Ba người chúng tôi đồng loạt nhìn Nene-chi với ánh mắt ``hết nói nổi.'' Kana-chan nhanh chóng chỉnh lại: ``Nó là mấy cái tiết khí đấy! Cứ tra tìm đi là biết!''

Nene-chi xị mặt xuống, có vẻ không hào hứng với những khái niệm học thuật này cho lắm. Nhận ra điều đó, Flare-chan liền hỏi một câu dễ hơn để kéo sự chú ý của em ấy trở lại: ``Mà nhắc đến mùa xuân, mọi người hay nghĩ đến gì?''

Chưa kịp để ai trả lời, ``ma mới'' đã lập tức bật dậy như lò xo và tuyên bố một cách hùng hồn:  ``Nene biết nhiều thứ về mùa xuân lắm chứ! Nhiều là cái chắc luôn!''

Câu nói tự tin quá mức ấy khiến Kanata-chan lập tức nảy ra ý định trêu chọc cô nàng một chút. Em ấy nhoẻn miệng hỏi: ``Vậy em có biết trong bài đồng dao Humpty Dumpty có đoạn nào nói về mùa xuân không?''

Nàng Hải cẩu chớp mắt, có vẻ hơi hoang mang. Thấy vậy, nàng Dị giới liền đọc hai câu ấy ra:

\begin{center}
  ``Mùa xuân cây cối xanh màu\\
  Ta toan nói rõ đôi câu tỏ bày\footnote{Nguyên văn: ``In Spring, when woods are getting green, I'll try and tell you what I mean.'' Đây thực chất là hai câu thơ của nhân vật Humpty Dumpty từ tác phẩm \textit{Through the Looking-Glass, and What Alice Found There} viết bởi Lewis Carroll.}.''
\end{center}

Tôi đổ mồ hôi hột, nhìn sang Nene-chi với vẻ lo lắng: ``Anh không nghĩ là em ấy sẽ hiểu hai câu đồng dao đó đâu\dots''

Đúng như dự đoán, Nene nheo mắt suy nghĩ một lúc, rồi sau đó đọc lại một cách vô cùng\dots\ sáng tạo khiến cho cả ba bọn tôi đổ mồ hôi hột:
\begin{center}
  ``Mùa xuân tới, khi\dots\ mừng?\\Ta ton nói rũ hai cầu tày bỏ!''
\end{center}

Đã thế lại còn đọc một cách vô tư nữa chứ. Chậc.

\img{8-10b}

Cả ba người chúng tôi lặng người trong vài giây, rồi đồng loạt đặt những dấu hỏi to đùng về khả năng nghe hiểu của em ấy. Chúng tôi đồng thanh: ``Em đang đọc cái quái gì thế?''

Thấy rằng Nene-chi có vẻ lươn lẹo quá mức, Kanata-chan quyết định chốt hạ một câu đầy tính khiêu khích: ``Thế tóm lại là em không biết một tí gì về mùa xuân chứ gì?''

Dĩ nhiên, nàng Hải cẩu không thể chịu thua dễ dàng như vậy. Em ấy lập tức phản bác: ``Biết chứ! Em còn nhìn thấy tận mắt cơ mà! Em thề!''

Tuy nhiên, có vẻ vì quá vội vã muốn chứng minh bản thân, nên em ấy hơi bị\dots\ thái quá, tới mức rơi nước mắt. Và khi cảm giác rằng không ai tin mình, cô nàng bỗng hậm hực quay đầu bỏ chạy.

\img{8-10c}

Flare-chan ngạc nhiên gọi với theo: ``\textbf{Nene-chi!?}''

Kanata cũng tỏ ra bối rối: ``Em ấy chạy đi đâu rồi kìa!?''

Còn tôi, không nghĩ nhiều, liền lập tức đuổi theo.

Hai người con gái phía sau còn chưa kịp phản ứng gì thì tôi đã lao đi mất. Nhưng cũng phải thôi. Bởi vì đây là thói quen của tôi rồi.

Nàng Hải cẩu vừa chạy vừa khóc, nước mắt lăn dài trên má, nức nở như một đứa trẻ đang bị oan ức. Em ấy vừa chạy vừa thốt lên: ``Mình biết về mùa xuân mà\dots! Sao không ai tin mình hết vậy!?''

Tôi thì không thể đứng yên được, vì nếu để em ấy chạy lung tung, không biết chuyện gì sẽ xảy ra. Mặc cho hơi thở ngày càng dồn dập, tôi vẫn cố gắng bám theo, trong khi cố gắng vọng lại: ``\textbf{Nene-chi! Dừng lại đi, chạy kiểu này nguy hiểm lắm!}''

Nhưng em ấy không nghe, chỉ cắm đầu chạy, mặc cho con đường phía trước ngày càng mờ dần.

\ct

Bỗng nhiên, cảnh vật xung quanh thay đổi. Không còn văn phòng, không còn lối đi quen thuộc nữa, chỉ có một không gian trắng xóa bao trùm lấy tất cả.

``Hả? Cái quái gì thế này!? Đây là đâu?''

Tôi kịp bắt lấy Nene-chi ngay khi em ấy bất ngờ ngã sấp mặt xuống sàn - hoặc ít nhất là thứ gì đó giống như sàn. Nhưng trước khi tôi kịp gọi, em ấy đã lồm cồm bò dậy, mắt vẫn còn đỏ hoe vì khóc.

Trong những cơn sụt sịt, dường như em ấy cũng định thần lại được, để rồi cũng giống tôi - ngơ ngác không biết đây là đâu.

Vừa lau nước mắt, vừa ngó nghiêng xung quanh, Nene-chi hỏi tôi: ``Ủa? Em vừa làm gì vậy?''

Tôi nhún vai đáp lại: ``Em vừa chạy đi đấy, anh ở đằng sau đuổi theo em thôi.''

Nene chớp chớp mắt, nhìn quanh không gian vô định này: ``Đây là đâu đây?''

``Chịu. Có giời mới biết.''

\img{8-10d}

Trong lúc còn tôi vẫn còn đang hoang mang không biết đây là nơi nào, thì ngay lúc này một giọng nói quen thuộc vang lên giữa khoảng không: ``Nene\dots\ Nene\dots''

Có vẻ như Giời đã nghe thấy lời thỉnh cầu của tôi. Có lẽ vậy.

Cả hai chúng tôi đồng loạt quay về hướng phát ra âm thanh ấy. Một bóng người dần hiện lên từ hư vô.

``Đây là vùng đất mùa xuân, và ta là nữ thần mùa xuân đây.''

Giọng nói này nghe quen quá\dots\ Tôi và Nene-chi chỉ tay thẳng vào người đó, không hẹn mà cùng đồng thanh: ``À! Miko-senpai/-chan kìa.''

Cái bóng lập tức bật dậy, hét lên đầy tức tối: ``\textbf{Này! Đừng có nói toẹt ra như vậy chứ, ne!}''

Tôi nhún vai đáp lại: ``Thì giọng bẹt của em nghe dễ đoán mà. Giống như hồi em làm show demo \textit{Ai là triệu phú} cho Sui-chan\footnote{Kuon đang nói đến buổi mini-live 3D của Suisei, ngày 21/8/2021.} ấy.''

Ngay sau đó, tôi nghe thấy một tiếng đập bàn cực mạnh, kèm theo tiếng rít khó chịu của míc-rô vọng lại. Có vẻ như em ấy đã bị tôi bắt bài.

Tôi vội vàng bịt tai lại, còn đồng đội của tôi thì bật cười khúc khích. Miko-chi đúng là dễ bị trêu thật.

Sau một hồi trấn tĩnh, Miko- à không, ``nữ thần mùa xuân'' nghiêm giọng tuyên bố, vẫn với giọng bẹt và âm ``ne'' ấy: ``Vì nhà ngươi có niềm tin vào mùa xuân, ta sẽ ban cho ngươi sức mạnh của mùa xuân!''

Nene-chi chớp mắt đầy ngạc nhiên, còn tôi thì cau mày: ``Khoan, sức mạnh mùa xuân là cái gì-''

Chưa kịp nói hết câu, cả người cô nàng Hải cẩu đột nhiên tỏa sáng với một thứ ánh sáng màu hồng rực rỡ. Em ấy ngạc nhiên nhìn hai tay mình phát sáng, còn tôi chỉ có thể đứng đó với một dấu chấm hỏi to đùng trên đầu rẳng thế tại sao lại là em ấy.

\img{8-10e}

``Nữ thần'' kia để lại lời nhắn cuối cùng trước khi biến mất: ``Hãy cho họ biết thế nào là lễ độ.''

Nene-chan siết tay lại đầy quyết tâm, mắt sáng lên: ``Em sẽ cố gắng ạ, Miko-senpai!''

Nhưng ngay lập tức, giọng nói của nữ thần vang lên đầy tuyệt vọng: ``\textbf{Đã bảo rồi! Ta là nữ thần mùa xuân cơ mà, ne!}''

Tôi bật cười, Nene cũng cười theo. Chúng tôi cười một tràng dài vì cái cách Miko-chi phản ứng khi bị chính chúng tôi bắt tại trận quá dễ thương.

Rồi, một ánh sáng màu hồng lại bao trùm lấy cả hai chúng tôi, những dải hoa anh đào rơi xuống, quay một vòng bao quanh chúng tôi.

\ct

Khi tôi nhận thức lại thì\dots\ chúng tôi đã trở về văn phòng.

Tôi được đưa trở về trước, vừa kịp lấy lại ý thức thì đã nghe thấy hai giọng nói đồng thanh vang lên: ``A, Kuon-senpai về rồi kìa!''

Flare-chan và Kana-chan cùng lúc thốt lên, cùng tỏ ra ngạc nhiên khi thấy tôi đột nhiên xuất hiện giữa phòng.

``Anh vừa đi đâu thế?'' nàng Dị giới hỏi.

Tôi chỉ đáp gọn: ``Đuổi theo Nene-chan. Tìm thấy em ấy rồi.''

Nàng Thiên sứ định hỏi tôi em ấy ở đâu, nhưng chưa kịp cất lời, một luồng ánh sáng hồng chói lóa lại bùng lên ngay phía sau tôi, và một giọng nói đầy kiêu hãnh vang lên: ``Hãy nhìn Super-Nene ta đây, giờ với sức mạnh mùa xuân mới!''

Nene xuất hiện với một tư thế cực ngầu - ít nhất là trong suy nghĩ của em ấy: Nene-chi từ đâu xuất hiện sau luồng ánh sáng hồng hào kia, khẽ lơ lửng trên không rồi nhẹ nhàng chạm đất. Tôi nhíu mày nhìn cô nàng, trong đầu đầy dấu chấm hỏi.

``Rồi, vậy sức mạnh mới đó là gì?''

Nhưng chưa kịp nghe câu trả lời, tôi đã thấy Nene đột nhiên tạo dáng, giơ hai tay lên và bắn ra một thứ tia sáng màu hồng thẳng vào đầu Thiên sứ.

\img{8-10f}

Tôi vội vàng chuẩn bị tinh thần để đỡ một trận la hét vì đau đớn\dots\ nhưng trái ngược với dự đoán, Kana-chan lại tròn mắt, tỏ ra bất ngờ: ``A! Chị hết đau đầu rồi này!''

Tôi chớp mắt ngạc nhiên. Hả? C-Cái gì?

Nene-chi cũng không kém phần kinh ngạc, em ấy nhìn bàn tay mình rồi nhìn nàng Thiên sứ với vẻ không tin nổi.

Tôi sửng sốt tới mức toát mồ hôi hột: ``V-Vậy cũng được nữa?''

Ngay lúc đó, một cơn gió nhẹ bất ngờ lướt qua căn phòng. Nhưng cơn gió này\dots\ lại đến từ một thực thể có hình dạng của một cánh hoa đào, lại còn được vẽ bằng giấy!? Nó trôi lơ lửng một cách đầy ma mị.

Flare liền hít một hơi sâu, đôi mắt ánh lên vẻ xúc động: ``Cơn gió xuân đầu tiên kìa!''

Tôi khoanh tay lại, cau mày nhìn cánh hoa đào giấy màu hồng kia: ``Sao anh cảm thấy có gì đó sai sai nhỉ?''

Nene-chi không trả lời, chỉ búng tay một cái, khiến cho cơn gió kỳ lạ ấy bay đi mất. Tôi bắt đầu cảm thấy tình huống này ngày càng kỳ quặc.

Thế rồi, tôi chưa kịp suy nghĩ thêm về tình huống ky quặc đó thì đã nghe thấy hai người con gái kia đồng loạt lên tiếng bằng một giọng điệu đầy ngưỡng mộ:

\begin{dialogue}
  \speak{Kanata} Nene-chan, cảm ơn em. Em là vị cứu tinh đấy.
  \speak{Flare} Một nữ thần! Là một nữ thần mùa xuân đó!
\end{dialogue}

\img{8-10g}

Tôi chỉ biết lắc đầu nhìn tình huống khó hiểu này, buột miệng: ``\dots Quá sai luôn.''

Nhưng, Kana-chan và Flare-chan không quan tâm đến lời của tôi. Họ bắt đầu hô vang: ``Nữ thần! Nữ thần!'' như là một đấng cứu tinh của thế giới.

Ban đầu, Nene-chi chỉ đứng đó, vẻ mặt ngạc nhiên như chưa kịp hiểu chuyện gì đang diễn ra. Nhưng rồi, chỉ sau vài giây, em ấy chớp mắt một cái, và lập tức lật mặt sang bản chất thật: một cô nàng kiêu căng chính hiệu.

Nene khoanh tay trước ngực, cười khoái chí: ``Vậy thì quỳ xuống và tôn thờ em đi! Ha!''

Flare và Kanata không chút do dự mà quỳ xuống, hai tay chắp lại như đang cầu nguyện.

Tôi chỉ có thể đứng đó, nhún vai, thở một hơi dài đầy bất lực.

``\textit{Mình không có tuổi để đọ với em ấy rồi\dots}''

Rồi, tôi khoanh tay lại, lắc đầu lẩm bẩm: ``Rồi chẳng biết `Super-Nene' này làm được trò trống gì không đây\dots''

%==========================================TRUYỆN PHỤ=========================================
\omk

Sau khi hoàn thành bản báo cáo và nộp cho A-san cùng các đồng nghiệp, tôi bước vào văn phòng, thở phào nhẹ nhõm. Ngẩng đầu lên, tôi thấy một thành viên mới của \hll\ đang say giấc ngon lành trên bàn làm việc.

Tôi nhìn gương mặt đang ngủ kia, khẽ bật cười: ``Ngủ trông cũng ngon nhỉ.''

Tiến lại gần hơn, tôi để ý thấy em ấy đang ôm một cuốn sách, trang đang mở đúng đoạn ``Bài đồng dao của Humpty Dumpty''.

\img{8-10h}

Tôi liếc nhìn tựa sách, nhíu mày. Quả thật, mùa xuân là mùa của văn chương. Nhất là với một thành viên tài năng như Nene-chi.

Tôi nghĩ thầm một lúc lâu, rồi tiếp tục đọc những dòng chữ mà em ấy đang giở.

Một người thích hội họa như em ấy, có vẻ cũng rất quan tâm đến thi ca. Nhưng rồi, giữa không gian yên tĩnh, tôi bất ngờ nghe thấy giọng ngái ngủ của em ấy lẩm bẩm đọc thơ: ``Mùa xuân đến, mừng\dots''

Tôi cứng đờ, đổ mồ hôi hột. Rõ ràng đây không phải là Humpty Dumpty.

Chậc. Có vẻ như em ấy còn phải học nhiều về thế giới này lắm, Momomomo-chan ạ. Và đúng, tôi cố tình đọc sai họ em ấy đấy, vì ngay từ đầu tôi đã nói: họ của Nene-chi rất dễ gây nhầm lẫn.

Tôi khẽ lắc đầu, đặt điện thoại xuống và ngồi đối diện em ấy, chờ em ấy tỉnh giấc.

\begin{center}
  \uline{40 phút sau.}
\end{center}

Cô nàng ngủ cùng cuốn sách từ từ tỉnh dậy. Dụi mắt vài cái, em ấy nhìn tôi đang ngồi đọc báo, rồi ngáp dài: ``Chào buổi sáng, Kuon-senpai\dots''

Tôi liếc nhìn đồng hồ, rồi nhìn em ấy với vẻ muốn chọc ghẹo em ấy đôi chút. ``Quá trưa rồi em ạ. Nhưng mà, chào buổi sáng, Suzumomo-chan.''

Nghe thấy tên họ của mình bị đọc sai, em ấy lập tức phản ứng, mặt nhăn lại: ``Em là Momosuzu!''

Tôi phì cười, lắc đầu: ``Xin lỗi, xin lỗi, anh biết rồi mà\dots\ Momotetsu-chan\~{}''

``\textbf{Mồ! Kuon-senpai!} Đừng có trêu em nữa!'' em ấy phồng má, rõ ràng đang dỗi.

Nhìn phản ứng đáng yêu của em ấy mỗi lần tôi cố tình đọc sai họ, tôi không nhịn được cười. Sau đó, tôi nghiêm túc lại một chút, giọng điệu trở nên nhẹ nhàng hơn: ``Thôi, đùa tí cho vui. Ngủ ngon không?''

Em ấy ngáp thêm lần nữa, rồi dụi mắt: ``Ngon lắm\dots''

``Ừm-hứm. À mà, sách gì vậy?'' tôi chỉ vào cuốn sách đang mở kia, cùng với những cuốn khác.

Em ấy ngước về một trong những cuốn sách đó, vừa ngáp vừa hớn hở giải thích: ``À, đây là cuốn tập truyện mà em mượn từ Flare-senpai á! Mùa xuân là mùa đọc sách mà nhỉ!''

Tôi phì cười: ``Anh lại thấy đây là mùa ngủ thì hợp lý hơn đấy.''

``Anh cứ đùa\dots'' Em ấy lại phồng má, có vẻ không đồng tình.

Tôi tựa người vào ghế, khoanh tay: ``Thế ngoài ra còn có gì nữa? Trông em ngủ ngon thế kia, chắc là mơ giấc mơ đẹp hử?''

Đôi mắt em ấy sáng lên khi nhắc về giấc mơ. ``Nene thấy tuyệt lắm! Em mơ thấy Kanata-senpai với Flare-senpai cùng em nói chuyện về mùa xuân cơ. Em còn mơ thấy mình được Miko-senpai ban phát sức mạnh mùa xuân nữa!''

Tôi nheo mắt, nghiêng đầu: ``Sao lại có cả Miko-chan ở đó?''

``Chị ấy tự nhận mình là thần mùa xuân mà!''

Tôi im lặng vài giây, sau đó gật gù. ``\dots Nghe cũng hợp lý. Họ của em ấy nghĩa là anh đào mà nhỉ.''

Trong lúc chúng tôi nói chuyện, ngay lúc đó, cánh cửa văn phòng bật mở, và Miko bước vào.

Tôi lập tức nhận ra ngay điểm kỳ lạ: trên đầu em ấy có đội một cái băng với dòng chữ ``CHÚA'' to đùng gắn trên một tờ giấy. Bộ trang phục quái dị đó là sao vậy?

Tôi còn chưa kịp phản ứng thì\dots

*VÚT!*

Một bóng dáng màu vàng với đôi mắt sáng như sao bỗng lao vụt qua như tia chớp, giật phăng chiếc băng khỏi đầu Miko-chi.

``\textbf{Này! Trả chị đây!!}'' nàng Vu nữ thét lên trong khi kẻ cắp đội ngay chiếc băng vào đầu mình.

``Oa! Trông mình cũng oách ghê nhỉ!''

\img{8-10i}

Không ai khác, chính là ``Hải cẩu'' Momosuzu Nene.

Tôi lặng người nhìn cảnh tượng trước mắt. Mikochan thì giơ tay cố chộp lại chiếc băng đầu, nhưng Nene-chi nhanh chóng né sang một bên, cười khoái chí.

``Nene! Nghiêm túc đi! Trả đây mau!!''

``Hahaha! Đố chị bắt được em đây nè\~{}''

Và thế là cuộc rượt đuổi bắt đầu. Người chị thì tuyệt vọng đòi lại chiếc mũ, còn người em thì thích thú khi trêu chọc tiền bối của mình.

Tôi khoanh tay đứng nhìn, thở dài: ``Thế này mà bảo là 'thần mùa xuân' đấy à\dots?''

Cảnh tượng này làm tôi nhớ lại một chuyện tương tự năm ngoái, lúc đó, Choco-sensei cũng từng bị Korone-chan trấn mất cái mũ, trên chiếc xe Lamb(o)rghini của Okayu-chan\footnote{Xem lại phần {\abv Truyện phụ}, {\flaregothic ÉPISODE 59}, quyển số Năm.}.

Đúng là truyền thống \hll\ mỗi khi có thành viên mới lên văn phòng mà.

Nhìn lại hai cô nàng đang rượt nhau quanh văn phòng, tôi lắc đầu, cười nhẹ. NePoLaBo đúng là không bao giờ làm tôi thất vọng.

Cả bốn người họ đều đã đặt chân vào công ty, mỗi người mang theo một sắc thái riêng biệt - một tính cách không lẫn vào đâu được. Nhưng dù khác biệt đến mấy, họ vẫn có chung một điểm: tất cả đều là những kẻ gây náo loạn theo cách riêng của mình.

``Haa\dots\ Thế là xong một chặng đường nữa.'' tôi lẩm bẩm, nhưng ngay sau đó, một suy nghĩ khác lại len lỏi vào tâm trí - một suy nghĩ có phần đáng sợ hơn nhiều, dù tôi đã luôn nghĩ về nó.

Bắt đầu từ hôm nay, văn phòng này sẽ còn hỗn loạn hơn trước, một cách không thể đoán trước được. Những tiếng cười, những trò quậy phá, những tình huống chẳng ai ngờ tới sẽ chẳng bao giờ ngừng lại. Và với sự gia nhập của họ, \hll\ càng lúc càng trở nên rực rỡ hơn, theo cách riêng của nó.

Nhưng NePoLaBo chưa bao giờ chỉ có bốn người. Dù không còn đứng trên sân khấu, dù đã bước đi trên một con đường khác, Aloe-chan vẫn sẽ luôn là một phần của họ. Và với tôi, em ấy vẫn luôn là một phần của gia đình nhà Hikawa này. Bốn người ở đây, một người ở nhà, vậy chẳng phải họ vẫn luôn là năm thành viên hay sao?

Tôi nhắm mắt, thở ra một hơi thật dài trước khi mở mắt trở lại. Những thế hệ tiếp theo rồi sẽ đến, từng người một. Hết nhóm này đến nhóm khác, mỗi người lại đem đến một sắc màu mới, một câu chuyện mới.

Tôi mỉm cười.

Phải, một nụ cười. Không phải một tiếng thở dài đầy chán nản, mà là một nụ cười thật sự. Vì suy cho cùng, tôi đã quen với chuyện này từ lâu rồi. Và nếu có một điều chắc chắn\dots\

\dots thì đó là: từ nay, cuộc sống văn phòng của tôi sẽ càng nhộn nhịp và hỗn loạn hơn.

Chẳng bao lâu nữa, mùa hè sẽ đến. Còn hơn hai tuần nữa thôi.

Chúng tôi cần phải bắt đầu chuẩn bị cho những hoạt động mùa hè sắp tới.

Nếu chúng tôi có thể.

\end{document}